\documentclass[11pt,letterpaper]{article}
\usepackage[T1]{fontenc}
\usepackage{lmodern,microtype}
\usepackage[margin=1in]{geometry}
\usepackage{amsmath,amssymb,amsthm}
\usepackage{booktabs,tabularx}
\usepackage{xcolor}
\usepackage[colorlinks=true,linkcolor=blue!40!black,citecolor=blue!40!black,
            urlcolor=blue!40!black]{hyperref}

\newcommand{\manuscripttitle}{Optimal pooled measurements with unknown values\\over abelian groups}
\newcommand{\manuscriptkeywords}{detecting matrix, coin weighing, abelian group, adversarial error, sparse recovery}
\hypersetup{
  pdftitle={Optimal pooled measurements with unknown values over abelian groups},
  pdfauthor={Joseph M. Shunia},
  pdfsubject={Detecting matrices for signals over abelian groups},
  pdfkeywords={\manuscriptkeywords}
}
\newtheorem{theorem}{Theorem}[section]
\newtheorem{lemma}[theorem]{Lemma}
\newtheorem{proposition}[theorem]{Proposition}
\newtheorem{corollary}[theorem]{Corollary}
\numberwithin{equation}{section}
\setlength{\emergencystretch}{1.5em}
\addtolength{\skip\footins}{6pt}

\title{\vspace{-1.5\baselineskip}%
  \manuscripttitle}
\author{Joseph M. Shunia\thanks{Independent researcher, Ann Arbor, Michigan, United States.
  Email: \href{mailto:jshunia@gmail.com}{\texttt{jshunia@gmail.com}}.}}
\date{November 22, 2026}

\begin{document}
\maketitle
\vspace{-\baselineskip}

\begin{abstract}
We study the number of subset sums needed to recover an unknown vector over an abelian group.
The vector has at most $s$ nonzero entries and at most $q$ distinct nonzero values, and the
values themselves are unknown. For each dimension and each substitution fraction below one
quarter, we construct by a probabilistic argument a single ordered binary matrix whose
prefixes attain the optimal order for all groups and all $s,q$. The order depends on the group
only through the least prime order of one of its elements. The same matrix gives signed measurements on
groups without an element of order two. For any specified matrix we prove an exact
cyclic-subgroup criterion for collisions and give a finite certificate by Smith normal form.
\\[2mm]
{\bf 2020 Mathematics Subject Classification:} 94B25 (primary), 05D40, 15B36 (secondary).
\\[2mm]
{\bf Keywords:} \manuscriptkeywords.
\end{abstract}

\section{Introduction}\label{sec:intro}

Four bags each contain either one counter or four counters. Measure each triple of bags. If
the instrument reports totals modulo nine, four ones and four fours both produce the readings
$(3,3,3,3)$. The same four tests determine every real-valued input. Which arithmetic features
create ambiguities like this, and how many tests suffice when the values in the bags are
themselves unknown?

Write $G$ additively and assume $G\ne 0$. Let $X_{s,q}(G)$ consist of vectors in $G^n$ with at
most $s$ nonzero coordinates and at most $q$ distinct nonzero values, where $1\le q\le s\le n$
and $n\ge2$. Different candidate vectors may use entirely different values. A subset test adds
the coordinates in a specified subset. A signed test gives every coordinate coefficient $+1$
or $-1$. A correct reading reports one exact element of $G$. A schedule's minimum distance is
the least Hamming distance between the answer words of distinct inputs.

Let $p(G)$ be the least prime $p$ for which $G$ has an element of order $p$, with
$p(G)=\infty$ for torsion-free groups. All logarithms are natural. With $\min(u,\infty)=u$,
define
\[
F(n,s,q,P)=
\frac{s\bigl(\log(en/s)+\log(1+\min(q,P))\bigr)}
     {\log(1+\min(s,P))}.
\]

\begin{theorem}\label{thm:rate}

Fix $0\le\eta<1/4$. For each $n$ there is a single ordered array of binary rows such that, for
every $s,q,G$, its first $\lceil C_\eta F(n,s,q,p(G))\rceil$ subset tests recover every member
of $X_{s,q}(G)$ through $\lfloor\eta m\rfloor$ substitutions, where $m$ is that prefix length.
The same binary array, with entries changed by $a\mapsto2a-1$, has the signed-test guarantee
on every group with no element of order two. A group with an element of order two admits
indistinguishable one-sparse inputs for every signed schedule. In both recoverable models even
a deterministic adaptive strategy with truthful answers needs at least $cF(n,s,q,p(G))$ tests
in the worst case, for an absolute $c>0$. One quarter is the limiting substitution fraction
for nonadaptive schedules as $n$ grows.
\end{theorem}

Detecting matrices encode a vector by sums of selected entries. Lindstr\"om studied the case
of a known finite alphabet \cite{lindstrom}. When real coin weights can take a bounded number
of unknown values, Alon and Kozlov's balanced weighing theorem already implies the optimal
order for dense noiseless recovery \cite{alon}. Sparse parity checks over a prescribed prime
field and noisy recovery of arbitrary real weights were studied by Bshouty and Mazzawi
\cite{bm,bmnoise}. For binary counterfeit weights Bshouty obtained the limiting one-quarter substitution threshold \cite{bnoise}. Here both the alphabet and the relations among its elements can vary. A
single integer matrix must control a possible collision in every characteristic while
retaining the gain from repeated values.

For real weights, $p(G)=\infty$. With $s=n$ and a fixed number of unknown values, the optimal
order is $n/\log n$. With $s=n$ and $q$ growing like $\log n$, it is $n\log\log n/\log n$.
Allowing a different value at every position brings the cost up to order $n$. For a counter
taking values modulo two, subset sums of dense inputs need order $n$ even if the only nonzero
value is one. The prefix length uses the assumed parameter bounds. The existence proof leaves
the computational cost of constructing a schedule and decoding its answers open.

To see the connection with balanced weighings more precisely, the difference between two dense
real inputs with $q$ nonzero values each takes at most $(q+1)^2$ values. A common offset makes
those values positive and preserves every balanced test. Alon and Kozlov's criterion forces
such a kernel vector to be constant, and a singleton test detects a nonzero constant. Their
inverse bound for $q<n^{1/4}$, with zero padding if necessary, gives $O(n\log(q+1)/\log n)$
subset tests. For $q\ge n^{1/4}$ the singleton tests already have this order. The issue for a
general group appears even in the four triple tests above. A collision modulo nine can survive
modulo three only after a change of labels.

\section{Cyclic witnesses and exceptional primes}\label{sec:cyclic}

A schedule corrects $t$ arbitrary substitutions precisely when its minimum distance exceeds
$2t$. Indeed, Hamming balls of radius $t$ about the possible transcripts must be disjoint. If
two transcripts differ in at most $2t$ positions, dividing those positions into two sets of
size at most $t$ produces a word in both balls.

The following fact holds for any family $X(G)\subseteq G^n$ preserved under coordinatewise
application of every zero-preserving set map $G\to K$ between the groups considered. Such a
map may merge several labels or erase them. Support restrictions, upper bounds on the number
of values and prescribed equalities between coordinates all have this property.

\begin{theorem}\label{thm:cyclic}

Fix an integer matrix $A$ and a set $E$ of its rows. Distinct admissible vectors agree on all
rows in $E$ over $G$ if and only if they do over at least one of its infinite cyclic or
prime-order cyclic subgroups. Consequently the minimum transcript distance is exactly the
minimum of the distances for those subgroups, identifying isomorphic cyclic groups. If a
family contains fewer than two vectors use $+\infty$ for its minimum distance.
\end{theorem}

\begin{proof}

Take a witnessing pair $x,y$ and put $z_i=x_i-y_i$. The subgroup $H$ generated by the $z_i$ is
nonzero and finitely generated. If it has a free summand, projection gives a nonzero
homomorphism $h:H\to\mathbb Z$, and $G$ contains an infinite cyclic subgroup. Otherwise $H$ is
finite and has a cyclic prime-order quotient $\mathbb F_p$, and $G$ contains a subgroup of
that order. Use such an $h$. At least one $h(z_i)$ is nonzero because the $z_i$ generate $H$.

The map is defined only on differences so far. For every coset of $H$ meeting the input labels
choose a representative $t$, with $t=0$ for $H$ itself. Relabel $a$ in that coset by
$f(a)=h(a-t)$ and extend arbitrarily to the remaining cosets, preserving zero. Each $x_i$ and
$y_i$ lie in the same coset. Thus $f(x_i)-f(y_i)=h(z_i)$. The new inputs are admissible,
differ somewhere, and still agree on every row of $E$ because $h$ is additive on $H$.
Conversely the subgroup embeds in $G$ and preserves all equalities and inequalities. Apply the
equivalence to the agreement rows of a minimum-distance pair. This proves the theorem.
\end{proof}

For the opening example, reducing 4 and 1 modulo three sends both to 1. Their differences
generate $\{0,3,6\}$ in $\mathbb Z/9$. Choosing 1 as the representative of their coset and
sending 3 to 1 sends the two labels to 1 and 0. Their triples still agree, now modulo three.
This illustrates why recentering the labels matters.

\subsection{A finite certificate for a given matrix}\label{sec:smith}

\begin{proposition}\label{prop:smith}

For fixed $n,s,q$, an integer matrix $A$, and a set $E$ of its rows, there are two
possibilities. Either $A_E$ identifies two distinct inputs over every nontrivial abelian
group, or there is a computable finite set of primes $\mathcal P$ such that a collision occurs
over $G$ exactly when $G$ contains an element of order $p$ for some $p\in\mathcal P$.
\end{proposition}

\begin{proof}

For $X_{s,q}$, prescribe the support and a label at each allowed position of each of the two
signals. There are finitely many such pairs of patterns. Use separate vertices for the unknown
labels of each signal and one common vertex for zero. Each coordinate is an oriented edge from
its label in $y$ to its label in $x$. Its value in $x-y$ is the difference of the endpoint
values. Choose a spanning forest of this graph and give its $d$ edges independent oriented
differences $w\in G^d$. Assigning these differences to the forest determines every edge
difference by adding along its unique forest path. Conversely every $w$ is realized by
assigning vertex potentials along the forest, starting at zero in each component and using the
zero vertex as root in its component. The patterns allow coincident or zero values, so every
assignment is legal.

In coordinates $x-y=Cw$ for an integer matrix $C$ with entries $0,1,-1$. Its forest rows form
an identity matrix. Ignore $d=0$, which yields only equal inputs. For a set $E$ of agreement
rows, form $D=A_E C$. Its integer Smith normal form has nonzero invariant factors
$a_1,\ldots,a_\rho$. Unimodular row and column operations are invertible on powers of every
abelian group. If $\rho<d$, choosing a nonzero element on a free column gives a collision in
every nontrivial group. If $\rho=d$, a nonzero kernel exists exactly when $G$ has an element
of prime order dividing some $a_i$. Indeed, $a_i g=0$ for some $g\ne0$ exactly when $G$
contains such a prime-order element. A nonzero kernel element gives distinct vectors because
the forest rows recover $w$. Taking the union over the finite pattern pairs gives all the
groups that fail these tests. Varying $E$ determines the full minimum-distance profile of the
schedule.
\end{proof}

For four tests each omitting one of four entries, the matrix is $J-I$, where $J$ is all ones.
Its invariant factors are $1,1,1,3$. It is injective even modulo two, yet modulo six inputs
differing by a constant two can be indistinguishable. Both of these groups have least prime
order two. The full set of prime orders determines the performance of an \emph{individual} schedule,
while the least one will suffice for the optimal order in Theorem~\ref{thm:rate}.

\section{Concentration with unknown labels}\label{sec:concentration}

\begin{lemma}\label{lem:point}

Work in $\mathbb Q$ or $\mathbb F_p$. Put $p=\infty$ in estimates for $\mathbb Q$. For signed
rows exclude $p=2$. Let $v$ have $r\ge1$ nonzero entries. For independent fair binary entries
or independent fair signs in a row $R$, uniformly over the target $a$,
\[
\Pr(Rv=a)\ \le\ Q_{r,p}:=\min\{1/2,\ 1/p+C/\sqrt r\}.
\]
\end{lemma}

\begin{proof}

Fixing all but one active entry gives the first bound. For a subset row over $\mathbb F_p$,
Fourier inversion and the triangle inequality bound the probability by
$p^{-1}\sum_{t=0}^{p-1}\prod_{j=1}^r|\cos(\pi t v_j/p)|$. Hölder's inequality with $r$ equal
exponents bounds this by $p^{-1}\sum_t|\cos(\pi t/p)|^r$, since multiplication by each $v_j$
permutes the residues and the absolute cosine is periodic in the residue. For signed rows use
$2v_j$ in place of $v_j$. The inequality $|\cos(\pi u)|\le\exp(-c\,\mathrm{dist}(u,\mathbb
Z)^2)$ gives a Gaussian sum bounded by $1/p+C/\sqrt r$ on comparing the positive terms with
the integral on $[0,\infty)$. Over $\mathbb Q$, clear denominators including that of the
target and reduce modulo arbitrarily large primes preserving all $r$ nonzero entries. Every
rational equality survives reduction. In particular $-\log Q_{r,p}\ge c_0 L_{r,p}$ for an
absolute $c_0>0$, where $L_{r,p}=\log(1+\min(r,p))$. This follows from the half bound at
bounded $\min(r,p)$ and the inverse square-root bound in that minimum for larger values.
\end{proof}

A probability estimate for one specified difference does not yet handle unspecified labels.
The following exposure argument chooses a witness after revealing a small collection of tests.
Throughout, $R_S$ denotes the restriction of a measurement row to the coordinate set $S$.

\begin{lemma}\label{lem:exposure}

Fix $\beta=1-2\eta>1/2$, a specified set $S$ of $r$ coordinates and an integer matrix $B$ with
$r$ rows. We consider witnesses $Bu$ having all $r$ coordinates nonzero. Suppose
$\operatorname{rank}B\le d$ in the field at hand. For $m$ independent rows, if
$d\le\epsilon_\eta m$ for a sufficiently small $\epsilon_\eta>0$, then
\[
\Pr\{\text{some such }Bu\text{ vanishes on at least }\beta m\text{ tests}\}
\le \exp(-c_\eta m L_{r,p}).
\]
\end{lemma}

\begin{proof}

To prove the bound, take a bad witness and select a basis from the rows $R_S B$ of its
vanishing tests. The basis uses an index set $I$ of size $h\le d$. For a fixed candidate $I$,
expose only its rows. If their joint kernel admits a $u$ with $Bu$ nonzero everywhere, choose
one by a fixed enumeration of the field vectors. Every row in the span of the exposed rows
kills this choice. Hence, whenever $I$ serves as the basis of a bad witness, at least $\beta
m-h$ of the remaining $m-h$ independent tests kill the chosen vector. This implication allows
us to test a single vector after exposure. It never conditions the remaining tests on $I$
actually being a basis.

Choose $\epsilon_\eta$ initially so small that
$(\beta-\epsilon_\eta)/(1-\epsilon_\eta)>1/2+\delta_\eta$ for some $\delta_\eta>0$. Each
remaining trial has probability at most $Q_{r,p}$. When $L_{r,p}$ is bounded, the binomial
upper tail above $1/2+\delta_\eta$ is at most $\exp(-c_\eta m)$ by the Chernoff bound. When
$L_{r,p}$ exceeds a constant depending only on $\eta$, the direct bound $2^m Q_{r,p}^{\,\beta
m-h}$ is at most $\exp(-c'_\eta m L_{r,p})$. The at most $\sum_{h\le d}\binom mh\le\exp(m
H(\epsilon_\eta))$ basis choices, with $H$ binary entropy in natural logarithms, are absorbed
by further decreasing $\epsilon_\eta$. Here $L_{r,p}\ge\log2$. This proves the estimate,
including rank zero and support one.
\end{proof}

\section{A single schedule for all parameters}\label{sec:prefix}

\begin{proof}[Proof of the upper bounds in Theorem~\ref{thm:rate}]

For the moment fix $s,q$ and an integer threshold $P\ge2$. We will cover $\mathbb Q$ and every
prime $p\ge P$ simultaneously. Signed rows cover the odd primes in that range. Only thresholds
up to $\max(2,s)$ will be needed. Put $F=F(n,s,q,P)$ and $m=\lceil K_\eta F\rceil$, where the
constant will be chosen large. Distinct inputs have a difference supported on some $S$ of size
$1\le r\le\min(2s,n)$.

If $q<P$, choose the two label assignments on $S$. There are at most $(q+1)^{2r}$ choices.
Their difference is $Bu$ with $2q$ unknown scalar labels in $u$ and at most one $+1$ and one
$-1$ in each row of $B$. A zero label contributes no variable. All coordinates outside $S$
cancel, so forgetting their labels only enlarges the event being bounded. Use $d=\min(2q,r)$.
If $q\ge P$, allow every difference on $S$ and take $B=I_r$, $d=r$. In the first case $F\ge q$
since $\log(1+t)/t$ decreases for $t>0$ and $q\le s$. In the second $F\ge s$. Thus $d/m$ can
be made arbitrarily small uniformly.

There are infinitely many primes but only finitely many distinguishable rank tests for one
pattern. For any agreement set $E$, the rows of $A_{E,S}B$, stacked with the rows of $B$, have
entries of absolute value at most $r$ and rank at most $d$ over $\mathbb Q$ and over every
prime field. Every minor of size at most $d$ is bounded in absolute value by $M=(r\sqrt d)^d$.
For $p>M$, every submatrix of this stack has the same rank modulo $p$ and over $\mathbb Q$, by
the minor criterion and Hadamard's inequality. If $Bu$ has full support modulo $p$ and the
rows in $E$ kill it, no row of $B$ belongs to the span of $A_{E,S}B$ modulo $p$. The same
holds over $\mathbb Q$. On the rational kernel each of the $r$ coordinate functionals is
therefore nonzero. Finitely many proper subspaces cannot cover a vector space over the
infinite field $\mathbb Q$, so a single rational witness has all $r$ entries nonzero. This
argument uses the same $E$ and creates no extra count of choices for $E$.

For a fixed support and pattern, union the probability estimate first over primes $P\le
p\le(r+1)^2$. There are at most $(r+1)^2$ of these and $L_{r,p}\ge L_{r,P}$. Each remaining
prime up to $M$, and $\mathbb Q$, has $L_{r,p}=\log(1+r)$ when $p>(r+1)^2$. There are at most
$1+M$ such cases. The logarithm of this count is $O(d\log(1+r))$ and is absorbed by increasing
$K_\eta$, since $F\ge d/2$. Primes exceeding both cutoffs are already covered by the rational
event. This division is important for small $P$, where paying the largest prime count against
the weakest probability bound would lose the claimed rate.

It remains to sum over the supports and patterns. Write $b=\min(q,P)$. Uniformly for $1\le r\le\min(2s,n)$,
\[
r\bigl(\log(en/r)+\log(1+b)\bigr)
\le C F L_{r,P}.
\]

For $\sqrt s\le r\le s$, the denominator $L_{r,P}$ is at least half of $L_{s,P}$ and
$r\log(en/r)$ is increasing on $[1,n]$. For $s\le r\le 2s$ use $L_{r,P}\ge L_{s,P}$ and
$r\le2s$. For $r<\sqrt s$, bound the numerator by $\sqrt s(\log(en/s)+\log s+\log(1+b))$ and
use $L_{r,P}\ge\log2$. Dividing this bound by $F$ costs at most a constant because
$(1+\log(1+s))^2/\sqrt s$ is bounded. This proves the comparison. It pays for $\binom nr$, the
label choices where used, and the factor $(r+1)^2$ above.

After increasing $K_\eta$, the failure probability for each $r$ is at most $\exp(-c_\eta
K_\eta F L_{r,P})$, with $c_\eta$ independent of the parameters. Also $F\ge c\log(n+1)$
uniformly. When $s\le\sqrt n$ this follows from $\log(en/s)$ and $s/\log(1+s)\ge1$. When
$s>\sqrt n$, use $F\ge s/\log(1+s)$ and $\sqrt n/\log(1+n)\ge c\log(1+n)$ with an absolute
$c>0$. Start with one ordered array of independent fair binary entries, using its
corresponding sign array for signed tests. The failure bounds concern its indicated prefixes
and require no independence between different prefixes. A union bound over at most $n$ support
sizes, $n^2$ pairs $(s,q)$, $\max(2,n)$ thresholds and the two row conventions gives total
probability less than one by choosing $K_\eta$ sufficiently large. Fix any successful array.

For a group $G$ use $P=\max(2,\min(s,p(G)))$, also when $p(G)=\infty$. This gives the same $F$
since $q\le s$. Its prefix handles the rationals and every required prime. Distances for
$\mathbb Z$ and $\mathbb Q$ agree by clearing denominators or by Theorem~\ref{thm:cyclic}.
Apply that theorem to any proposed pair with at least $(1-2\eta)m$ equal answers. The cyclic
prefix guarantees rule out every such pair. Its minimum distance thus exceeds $2\eta m$, which
suffices through $\lfloor\eta m\rfloor$ substitutions.
\end{proof}

\section{Lower bounds and the substitution threshold}\label{sec:lower}

\begin{proof}[Proof of the lower bounds in Theorem~\ref{thm:rate}]

Choose $g$ of order $p=p(G)$ if finite and otherwise of infinite order. Use letters
$0,g,\ldots,ug$ with $u=q$ in the infinite case and $u=\min(q,p-1)$ otherwise. The family
includes at least $(u+1)^s$ vectors on the first $s$ coordinates and at least $\binom ns$
binary support choices. Taking the larger shows that its logarithmic size is at least $c
s(\log(en/s)+\log(1+\min(q,p)))$. Indeed $(u+1)\ge2$, $\binom ns\ge(n/s)^s$ and $\log(u+1)$ is
comparable to $\log(1+\min(q,p))$. Each signed or subset query has at most $\min(p,2su+1)$
answers, omitting $p$ in the infinite case. Its log is at most $C\log(1+\min(s,p))$ since
$u\le s$. The same bound applies at every node of an adaptive truthful decision tree. Counting
its leaves proves the claimed lower bound.

For signed tests an order-two element $g$ makes $ge_1$ and $ge_2$ identical on every row.
Otherwise choose any $g\ne0$. On each signed row exactly one of the pairs $(ge_1,ge_2)$ and
$(ge_1,-ge_2)$ differs, so some pair is at distance at most $\lfloor m/2\rfloor$. For subset
tests, choose any $g\ne0$. The $n+1$ inputs $0,ge_1,\ldots,ge_n$ give binary answer words with
symbols $0,g$. In each row at most $(n+1)^2/4$ unordered pairs differ. Averaging gives a pair
at distance at most $m(n+1)/(2n)$. Distinguishing these binary words requires
$m\ge\log_2(n+1)$. Thus rounding vanishes as $n$ grows, proving the asymptotic upper limit
$1/4$ already at $s=q=1$. The same distance argument permits a mixture of substitutions and
known erasures whenever twice the substitution count plus the erasure count is less than the
guaranteed distance.
\end{proof}

\section*{AI Assistance Disclosure}

OpenAI's GPT-6.2 Astra was used in an autoresearch loop and reported this discovery to the author.

\begingroup\small
\begin{thebibliography}{9}
\bibitem{lindstrom}
B. Lindstr\"om, \emph{On a combinatorial problem in number theory}, Canadian Mathematical
Bulletin \textbf{8} (1965), 477--490.
\href{https://doi.org/10.4153/CMB-1965-034-2}{doi:10.4153/CMB-1965-034-2}.

\bibitem{alon}
N. Alon and D. N. Kozlov, \emph{Coins with arbitrary weights}, Journal of Algorithms
\textbf{25} (1997), 162--176.
\href{https://doi.org/10.1006/jagm.1997.0871}{doi:10.1006/jagm.1997.0871}.

\bibitem{bm}
N. H. Bshouty and H. Mazzawi, \emph{On parity check $(0,1)$-matrix over $\mathbb Z_p$}, SIAM
Journal on Discrete Mathematics \textbf{29} (2015), 631--657.
\href{https://doi.org/10.1137/120881129}{doi:10.1137/120881129}.

\bibitem{bmnoise}
N. H. Bshouty and H. Mazzawi, \emph{Algorithms for the coin weighing problems with the
presence of noise}, Electronic Colloquium on Computational Complexity, Report TR11-124 (2011).
\href{https://eccc.weizmann.ac.il/report/2011/124/}{ECCC TR11-124}.
\bibitem{bnoise}
N. H. Bshouty, \emph{On the coin weighing problem with the presence of noise}, in
APPROX/RANDOM 2012, Lecture Notes in Computer Science \textbf{7408}, 471--482.
\href{https://doi.org/10.1007/978-3-642-32512-0_40}{doi:10.1007/978-3-642-32512-0\_40}.
\end{thebibliography}
\endgroup

\end{document}
